% concatenated from abdelhai.tex
\documentclass[11pt]{article}

\usepackage[T1]{fontenc}
\usepackage{lmodern}
\usepackage{amsmath,amssymb,amsthm,mathtools}
\usepackage{geometry}
\usepackage{microtype}
\usepackage{booktabs}
\usepackage{enumitem}
\usepackage{hyperref}

\geometry{margin=1in}
\hypersetup{
  colorlinks=true,
  linkcolor=blue,
  citecolor=blue,
  urlcolor=blue
}

\newtheorem{theorem}{Theorem}[section]
\newtheorem{proposition}[theorem]{Proposition}
\newtheorem{lemma}[theorem]{Lemma}
\newtheorem{corollary}[theorem]{Corollary}
\theoremstyle{definition}
\newtheorem{definition}[theorem]{Definition}
\theoremstyle{remark}
\newtheorem{remark}[theorem]{Remark}
\newtheorem{example}[theorem]{Example}

\newcommand{\N}{\mathbb N}
\newcommand{\Z}{\mathbb Z}
\newcommand{\Q}{\mathbb Q}
\newcommand{\Int}{\operatorname{Int}}
\newcommand{\coeff}[2]{[x^{#1}]\,#2}

\title{\bfseries
Shifted-Binomial Expansions of Dilated Binomial Polynomials:\\
Multinomial Decimation, Reflection Symmetry, and Universal Divisibility}

\author{Abdelhai Doukali\\
\small Independent Researcher, Fez, Morocco\\
\small \texttt{doukaliabdelhai@gmail.com}}

\date{}

\begin{document}
\maketitle

\begin{abstract}
For integers $m\ge2$ and $r\ge1$, put $d=m+r-1$ and define the shifted-binomial
coordinates $B_{m,r,k}$ by
\[
 \binom{mn}{d}
 =
 \sum_{k=1}^{d+1}B_{m,r,k}\binom{n+k-1}{d}.
\]
The main purpose of this paper is to identify these coordinates with a
specific residue-class decimation of a multinomial coefficient sequence. We
prove
\[
 B_{m,r,k}
 =
 [x^{mk-1}](1+x+\cdots+x^{m-1})^{m+r}.
\]
This identity converts the alternating finite-difference formula for the
coordinates into a positive coefficient formula, and yields their exact
support. We then prove that the nonzero coefficient vector is palindromic
if and only if $r\equiv2\pmod m$. This symmetry criterion is derived here directly within the
multinomial-decimation framework. The principal arithmetic results are an unconditional divisibility law and
an exact formula for the greatest common divisor of each nonzero row. We prove
\[
 \frac{m}{\gcd(m,r-1)}\mid B_{m,r,k},
\]
and, more precisely, determine
\[
 \gcd_{1\le k\le K_{m,r}} B_{m,r,k}
\]
as an explicit product of prime powers determined by the $p$-adic valuations
of $m$ and $m+r-1$. Finally, a roots-of-unity filter gives the exact row sum
\[
 \sum_k B_{m,r,k}=m^{m+r-1}.
\]
Thus the shifted-binomial coordinates form a positive, arithmetically
structured decimation of an $m$-nomial coefficient sequence.
\end{abstract}

\noindent\textbf{Keywords:}
binomial coefficients; shifted-binomial basis; integer-valued polynomials;
multinomial coefficients; multisection; reflection symmetry; divisibility.

\medskip
\noindent\textbf{MSC 2020:}
05A10, 05A15, 05A19, 11B83, 13F20.

\section{Introduction}

Let
\[
 \Int(\Z)=\{P\in\Q[X]:P(\Z)\subseteq\Z\}
\]
denote the ring of integer-valued polynomials. The classical binomial
polynomials form a $\Z$-basis of this ring; see, for example,
\cite{Polya1919,Ostrowski1919,CahenChabert1997}. For a fixed degree $d$,
the shifted family
\[
 \left\{\binom{X+k-1}{d}:1\le k\le d+1\right\}
\]
is a $\Q$-basis of the vector space of polynomials of degree at most $d$.

We study the dilated binomial polynomial
\[
 g_{m,r}(X)=\binom{mX}{m+r-1},
 \qquad m\ge2,\quad r\ge1,
\]
and write
\[
 d=m+r-1.
\]
The corresponding shifted-binomial coordinates are defined by
\begin{equation}
 \label{eq:main-expansion}
 \binom{mn}{d}
 =
 \sum_{k=1}^{d+1}B_{m,r,k}\binom{n+k-1}{d}.
\end{equation}

The finite-difference inversion of this expansion gives an explicit
alternating sum for $B_{m,r,k}$. The main structural result of the present
paper is stronger conceptually: the same coefficient is a coefficient of a
power of the all-one polynomial,
\begin{equation}
 \label{eq:main-decimation}
 B_{m,r,k}
 =
 [x^{mk-1}](1+x+\cdots+x^{m-1})^{m+r}.
\end{equation}
Thus the coordinate vector in \eqref{eq:main-expansion} is exactly the
subsequence of the $(m+r)$-fold multinomial coefficient sequence in the
residue class $-1\pmod m$. The underlying multinomial coefficients are
classical; the contribution here is the precise identification of these
coefficients with the shifted-binomial coordinates of the dilated binomial
polynomial.

The multinomial-decimation framework also yields the reflection criterion
and, in addition, the stronger unconditional divisibility law
\begin{equation}
\label{eq:main-div}
\frac{m}{\gcd(m,r-1)}\mid B_{m,r,k}.
\end{equation}

The decimation identity also gives an exact description of the support.
If
\[
 K_{m,r}=d+1-\left\lceil\frac dm\right\rceil
 =m+r-1-\left\lceil\frac{r-1}{m}\right\rceil,
\]
then
\[
 B_{m,r,k}>0\quad\Longleftrightarrow\quad 1\le k\le K_{m,r}.
\]
The same representation yields the reflection criterion and, by a
roots-of-unity filter, the row-sum identity
\[
 \sum_{k=1}^{K_{m,r}}B_{m,r,k}=m^{m+r-1}.
\]

The paper is organized as follows. Section~\ref{sec:expansion} develops the
shifted-binomial expansion and the finite-difference formula. Section
\ref{sec:decimation} proves the multinomial-decimation identity and the exact
support. Section~\ref{sec:symmetry} proves the reflection criterion, including
the necessary strict-unimodality input. Section~\ref{sec:divisibility}
establishes both the universal divisibility law and the exact row-gcd formula.
Section~\ref{sec:rows} derives the row sums and records examples. We conclude
with a discussion of the scope of the results and possible extensions.

\section{Shifted-binomial coordinates}
\label{sec:expansion}

Throughout, $m\ge2$, $r\ge1$, and
\[
 d=m+r-1.
\]

\begin{lemma}
\label{lem:basis}
For every integer $d\ge0$, the polynomials
\[
 \binom{X+k-1}{d},\qquad 1\le k\le d+1,
\]
form a basis of the $\Q$-vector space of polynomials of degree at most $d$.
\end{lemma}

\begin{proof}
All members have degree $d$. Suppose
\[
 \sum_{k=1}^{d+1}c_k\binom{X+k-1}{d}=0.
\]
Evaluating successively at $X=0,1,\ldots,d$ gives a triangular system.
At $X=0$ only the term $k=d+1$ can be nonzero; hence $c_{d+1}=0$.
At $X=1$ the only remaining possible term at the diagonal is $k=d$,
so $c_d=0$, and induction gives $c_k=0$ for every $k$.
There are $d+1$ polynomials in a $(d+1)$-dimensional space.
\end{proof}

\begin{lemma}
\label{lem:gf}
For $1\le k\le d+1$,
\[
 \sum_{n\ge0}\binom{n+k-1}{d}x^n
 =
 \frac{x^{d+1-k}}{(1-x)^{d+1}}.
\]
\end{lemma}

\begin{proof}
Since
\[
 \sum_{j\ge0}\binom{j}{d}x^j
 =
 \frac{x^d}{(1-x)^{d+1}},
\]
and $\binom{j}{d}=0$ for $0\le j<d$, shifting the index gives
\[
 \sum_{n\ge0}\binom{n+k-1}{d}x^n
 =
 x^{-(k-1)}\frac{x^d}{(1-x)^{d+1}}.
\]
\end{proof}

\begin{theorem}[Finite-difference inversion]
\label{thm:inversion}
Let $f\in\Q[X]$ have degree at most $d$ and write
\[
 f(n)=\sum_{k=1}^{d+1}c_k\binom{n+k-1}{d}.
\]
Then
\[
 c_k
 =
 \sum_{j=0}^{d+1-k}
 (-1)^j\binom{d+1}{j}f(d+1-k-j).
\]
\end{theorem}

\begin{proof}
By Lemma~\ref{lem:gf},
\[
 \sum_{n\ge0}f(n)x^n
 =
 \frac{1}{(1-x)^{d+1}}
 \sum_{k=1}^{d+1}c_kx^{d+1-k}.
\]
Multiplication by $(1-x)^{d+1}$ and comparison of the coefficient of
$x^{d+1-k}$ gives the stated formula.
\end{proof}

\begin{corollary}
\label{cor:closed}
For $1\le k\le d+1$,
\begin{equation}
\label{eq:closed}
 B_{m,r,k}
 =
 \sum_{j=0}^{k-1}
 (-1)^j\binom{d+1}{j}
 \binom{m(k-j+1)+r-2}{d}.
\end{equation}
\end{corollary}

\begin{proof}
Apply Theorem~\ref{thm:inversion} to $f(n)=\binom{mn}{d}$. After the change of
index $j=d+1-k-\ell$, the terms with $\ell\ge k$ vanish, and the resulting
finite sum is exactly \eqref{eq:closed}.
\end{proof}

\begin{remark}
Formula \eqref{eq:closed} is an alternating finite-difference expression.
The multinomial-decimation theorem below replaces it by a nonnegative
coefficient formula, which is the main structural improvement of this
description.
\end{remark}

\section{Multinomial decimation and exact support}
\label{sec:decimation}

\begin{theorem}[Multinomial decimation]
\label{thm:decimation}
For every $1\le k\le d+1$,
\begin{equation}
\label{eq:dec}
 B_{m,r,k}
 =
 [x^{mk-1}](1+x+\cdots+x^{m-1})^{d+1}.
\end{equation}
In particular, $B_{m,r,k}\in\Z_{\ge0}$.
\end{theorem}

\begin{proof}
Use
\[
 1+x+\cdots+x^{m-1}=\frac{1-x^m}{1-x}.
\]
Then
\[
 (1+x+\cdots+x^{m-1})^{d+1}
 =
 (1-x^m)^{d+1}(1-x)^{-(d+1)}.
\]
Consequently,
\[
\begin{aligned}
 [x^{mk-1}](1+x+\cdots+x^{m-1})^{d+1}
 &=
 \sum_{j=0}^{k-1}
 (-1)^j\binom{d+1}{j}
 [x^{m(k-j)-1}](1-x)^{-(d+1)}\\
 &=
 \sum_{j=0}^{k-1}
 (-1)^j\binom{d+1}{j}
 \binom{m(k-j)+d-1}{d}.
\end{aligned}
\]
Since $d=m+r-1$,
\[
 m(k-j)+d-1=m(k-j+1)+r-2,
\]
so the last expression is exactly \eqref{eq:closed}.
Nonnegativity and integrality follow from the coefficient interpretation.
\end{proof}

\begin{remark}
Theorem~\ref{thm:decimation} identifies $B_{m,r,k}$ with the coefficient of
$x^{mk-1}$ in the $(m+r)$-th power of the $m$-nomial polynomial
$1+x+\cdots+x^{m-1}$. Thus the row is obtained by sampling one residue
class, namely $-1\pmod m$, of a classical multinomial coefficient sequence.
The novelty lies in this identification with the shifted-binomial
coordinates, rather than in the multinomial coefficients themselves.
\end{remark}

\begin{proposition}[Exact support]
\label{prop:support}
Put
\[
 K_{m,r}
 =
 d+1-\left\lceil\frac dm\right\rceil
 =
 m+r-1-\left\lceil\frac{r-1}{m}\right\rceil.
\]
Then
\[
 B_{m,r,k}>0\quad\Longleftrightarrow\quad 1\le k\le K_{m,r},
\]
and $B_{m,r,k}=0$ for $K_{m,r}<k\le d+1$.
\end{proposition}

\begin{proof}
Let
\[
 H(x)=(1+x+\cdots+x^{m-1})^{d+1}.
\]
Its degree is
\[
 D=(m-1)(d+1).
\]
Every coefficient of $H$ between degrees $0$ and $D$ is positive, because
each factor has positive coefficients in every degree from $0$ to $m-1$.
By Theorem~\ref{thm:decimation},
\[
 B_{m,r,k}>0
 \quad\Longleftrightarrow\quad
 0\le mk-1\le D.
\]
The lower inequality is automatic for $k\ge1$, while the upper inequality is
equivalent to
\[
 k\le d+1-\frac dm.
\]
Taking the largest integer satisfying this inequality gives
\[
 K_{m,r}=d+1-\left\lceil\frac dm\right\rceil.
\]
Substitution of $d=m+r-1$ gives the second expression.
\end{proof}

\section{Reflection symmetry}
\label{sec:symmetry}

\begin{definition}
The nonzero coefficient vector
\[
 (B_{m,r,1},\ldots,B_{m,r,K_{m,r}})
\]
is called palindromic if
\[
 B_{m,r,k}=B_{m,r,K_{m,r}+1-k}
\]
for every $1\le k\le K_{m,r}$.
\end{definition}

The converse direction of the symmetry theorem requires only a very local
monotonicity statement, not a general unimodality theorem.

\begin{lemma}
\label{lem:local-monotonicity}
Let
\[
 H(x)=(1+x+\cdots+x^{m-1})^N,\qquad N\ge1,
\]
and write $h_j=[x^j]H(x)$. Then
\[
 h_0<h_1<\cdots<h_{m-1}.
\]
\end{lemma}

\begin{proof}
For $0\le j\le m-1$, the upper bound $m-1$ on each exponent is
automatically satisfied when the total degree is $j$. Hence
\[
 h_j=[x^j](1-x)^{-N}=\binom{N+j-1}{j}.
\]
The latter sequence is strictly increasing in $j$ because
\[
 \frac{h_{j+1}}{h_j}=\frac{N+j}{j+1}>1
 \qquad (0\le j<m-1).
\]
\end{proof}

\begin{theorem}[Reflection criterion]
\label{thm:symmetry}
The vector
\[
 (B_{m,r,1},\ldots,B_{m,r,K_{m,r}})
\]
is palindromic if and only if
\[
 r\equiv2\pmod m.
\]
\end{theorem}

\begin{proof}
Let
\[
 H(x)=(1+x+\cdots+x^{m-1})^{m+r}
     =\sum_{j=0}^{D}h_jx^j,
 \qquad D=(m-1)(m+r).
\]
By Theorem~\ref{thm:decimation},
\[
 B_{m,r,k}=h_{mk-1}.
\]
Since $H$ is self-reciprocal,
\[
 h_j=h_{D-j}.
\]

Set
\[
 a_1=m-1,\qquad a_K=mK_{m,r}-1.
\]
The next sampled exponent is $a_K+m$, which is strictly larger than $D$.
Hence
\[
 s:=D-a_K
\]
is an integer satisfying
\[
 0\le s\le m-1.
\]
Moreover $a_1=m-1\le D/2$. Thus both $s$ and $a_1$ lie in the increasing
half of the symmetric coefficient sequence.

If the sampled vector is palindromic, then
\[
 h_{a_1}=h_{a_K}=h_{D-a_K}=h_s.
\]
By Lemma~\ref{lem:local-monotonicity}, this forces
\[
 s=m-1.
\]
Therefore
\[
 D-(mK_{m,r}-1)=m-1,
\]
or
\[
 D=m(K_{m,r}+1)-2.
\]
Reducing modulo $m$ gives
\[
 -r\equiv-2\pmod m,
\]
hence $r\equiv2\pmod m$.

Conversely, suppose $r\equiv2\pmod m$ and write
\[
 r-2=qm,\qquad q\ge0.
\]
Then
\[
 d=m+r-1=m(q+1)+1,
 \qquad
 K_{m,r}=d-q-1.
\]
The identity
\[
 D=m(K_{m,r}+1)-2
\]
follows directly. Hence
\[
 D-(mk-1)=m(K_{m,r}+1-k)-1.
\]
Using $h_j=h_{D-j}$,
\[
 B_{m,r,k}
 =h_{mk-1}
 =h_{m(K_{m,r}+1-k)-1}
 =B_{m,r,K_{m,r}+1-k}.
\]
Thus the vector is palindromic.
\end{proof}

\begin{remark}
The reflection criterion follows directly from the self-reciprocity of the
multinomial generating polynomial together with the local monotonicity
established above.
\end{remark}

\section{Universal divisibility}
\label{sec:divisibility}

The next theorem gives a universal divisibility law for all $m\ge2$ and
$r\ge1$.

\begin{lemma}
\label{lem:intz}
Let $P\in\Int(\Z)$ have degree at most $d$, and write
\[
 P(X)=\sum_{k=1}^{d+1}c_k\binom{X+k-1}{d}.
\]
Then every $c_k$ is an integer.
\end{lemma}

\begin{proof}
Apply Theorem~\ref{thm:inversion}. Each $c_k$ is an integer linear
combination of values of $P$ at integers, and these values are integers.
\end{proof}

\begin{theorem}[Universal divisibility]
\label{thm:div}
For every $m\ge2$, $r\ge1$, and $1\le k\le d+1$,
\[
 \frac{m}{\gcd(m,r-1)}\mid B_{m,r,k}.
\]
\end{theorem}

\begin{proof}
The identity
\[
 d\binom{mX}{d}
 =
 mX\binom{mX-1}{d-1}
\]
with $d=m+r-1$ suggests the integer-valued polynomial
\[
 Q(X)=X\binom{mX-1}{d-1}.
\]
For every $X\in\Z$, $Q(X)\in\Z$, so $Q\in\Int(\Z)$. Expand
\[
 Q(X)=\sum_{k=1}^{d+1}A_k\binom{X+k-1}{d}.
\]
By Lemma~\ref{lem:intz}, $A_k\in\Z$.

On the other hand,
\[
 d\,g_{m,r}(X)=mQ(X).
\]
Comparing the coordinates in the same shifted-binomial basis gives
\[
 dB_{m,r,k}=mA_k.
\]
Let
\[
 g=\gcd(m,d)=\gcd(m,r-1),
 \qquad
 m=gm_0,\quad d=gd_0.
\]
Then $\gcd(m_0,d_0)=1$ and
\[
 d_0B_{m,r,k}=m_0A_k.
\]
Euclid's lemma gives $d_0\mid A_k$. Write
\[
 A_k=d_0C_k,\qquad C_k\in\Z.
\]
Then
\[
 B_{m,r,k}=m_0C_k
 =\frac{m}{\gcd(m,r-1)}C_k,
\]
as required.
\end{proof}

\begin{remark}
Theorem~\ref{thm:div} gives a universal divisor, but it need not be sharp.
The exact greatest common divisor is obtained below by combining the
shifted-binomial basis with a fixed-divisor calculation for
$\binom{mX}{d}$.
\end{remark}

\begin{lemma}[Fixed divisor and shifted-binomial coordinates]
\label{lem:fixed-divisor}
Let $P\in\Int(\Z)$ have degree at most $d$, and write
\[
 P(X)=\sum_{k=1}^{d+1}c_k\binom{X+k-1}{d},
 \qquad c_k\in\Z.
\]
Then
\[
 \gcd_{1\le k\le d+1}c_k
 =
 \gcd_{n\in\Z}P(n).
\]
\end{lemma}

\begin{proof}
Let $G=\gcd_k c_k$ and let $\delta=\gcd_{n\in\Z}P(n)$.
Since each $\binom{n+k-1}{d}$ is an integer, the expansion of $P(n)$
shows that $G\mid P(n)$ for every $n$, and hence $G\mid\delta$.

Conversely, Theorem~\ref{thm:inversion} expresses each $c_k$ as an
integer linear combination of the values
\[
 P(0),P(1),\ldots,P(d+1-k).
\]
Thus $\delta\mid c_k$ for every $k$, so $\delta\mid G$. Therefore
$G=\delta$.
\end{proof}

\begin{lemma}[Minimum valuation of a product of consecutive integers]
\label{lem:consecutive-vp}
Let $p$ be prime and $L\ge1$. Then
\[
 \min_{t\in\Z}\sum_{j=0}^{L-1}v_p(t-j)=v_p(L!).
\]
\end{lemma}

\begin{proof}
For every $\ell\ge1$, among any $L$ consecutive integers there are at
least $\lfloor L/p^\ell\rfloor$ terms divisible by $p^\ell$. Hence
\[
 \sum_{j=0}^{L-1}v_p(t-j)
 \ge
 \sum_{\ell\ge1}\left\lfloor\frac{L}{p^\ell}\right\rfloor
 =v_p(L!).
\]
The lower bound is obtained by counting, for each $\ell\ge1$, the
terms that are divisible by $p^\ell$.

Equality is attained at $t=L$, because
\[
 \sum_{j=0}^{L-1}v_p(L-j)=v_p(L!).
\]
\end{proof}

\begin{theorem}[Exact row gcd]
\label{thm:exact-gcd}
Put
\[
 d=m+r-1,
 \qquad
 G_{m,r}:=\gcd_{1\le k\le K_{m,r}}B_{m,r,k}.
\]
For every prime $p\mid m$, set
\[
 a_p=v_p(m),
 \qquad
 q_p=\left\lfloor\frac{d-1}{p^{a_p}}\right\rfloor
     =\left\lceil\frac{d}{p^{a_p}}\right\rceil-1.
\]
Then
\begin{equation}
\label{eq:exact-gcd}
 G_{m,r}
 =
 \prod_{p\mid m}
 p^{\,a_p+v_p(q_p+1)-v_p(d)}.
\end{equation}
Equivalently,
\[
 G_{m,r}
 =
 \prod_{p\mid m}
 p^{
 v_p(m)
 +v_p\!\left(
 \left\lceil\frac{m+r-1}{p^{v_p(m)}}\right\rceil
 \right)
 -v_p(m+r-1)
 }.
\]
In particular,
\[
 \frac{m}{\gcd(m,r-1)}\mid G_{m,r}.
\]
\end{theorem}

\begin{proof}
By Proposition~\ref{prop:support}, the coefficients outside
$1\le k\le K_{m,r}$ vanish. Hence
\[
 G_{m,r}
 =
 \gcd_{1\le k\le d+1}B_{m,r,k}.
\]
Applying Lemma~\ref{lem:fixed-divisor} to
\[
 P(X)=\binom{mX}{d}
\]
gives
\begin{equation}
\label{eq:fixed-divisor}
 G_{m,r}
 =
 \gcd_{n\in\Z}\binom{mn}{d}.
\end{equation}
Thus it remains to compute the fixed divisor prime by prime.

Let $p$ be a prime and write
\[
 a=v_p(m).
\]
If $a=0$, multiplication by $m$ is a bijection on $\Z_p$. Therefore
\[
 \min_{n\in\Z_p}v_p\binom{mn}{d}
 =
 \min_{t\in\Z_p}v_p\binom{t}{d}.
\]
Taking $t=d$ gives value $1$, while more importantly the polynomial
$\binom{t}{d}$ has an integer value equal to $1$ at $t=d$. Hence the
minimum is $0$, and therefore
\[
 v_p(G_{m,r})=0
 \qquad(p\nmid m).
\]

Now suppose $a=v_p(m)>0$. Write
\[
 m=p^a u,\qquad p\nmid u,
\]
and put
\[
 q=\left\lfloor\frac{d-1}{p^a}\right\rfloor.
\]
For an integer $n$,
\[
 \binom{mn}{d}
 =
 \frac{\prod_{i=0}^{d-1}(mn-i)}{d!}.
\]
If $p^a\nmid i$, then $v_p(i)<a\le v_p(mn)$, and hence
\[
 v_p(mn-i)=v_p(i).
\]
For $i=p^aj$, with $0\le j\le q$,
\[
 mn-i=p^a(un-j),
\]
so
\[
 v_p(mn-i)=a+v_p(un-j).
\]
Consequently
\[
\begin{aligned}
 v_p\binom{mn}{d}
 &=
 \frac{}{}\sum_{\substack{1\le i<d\\p^a\nmid i}}v_p(i)
 +a(q+1)
 +\sum_{j=0}^{q}v_p(un-j)
 -v_p(d!).
\end{aligned}
\]
Since $u$ is coprime to $p$, multiplication by $u$ permutes the residue
classes modulo every power of $p$. Hence, for the purpose of minimizing
the finite valuation sum, $un$ may be replaced by an arbitrary integer
$t$. By Lemma~\ref{lem:consecutive-vp},
\[
 \min_n\sum_{j=0}^{q}v_p(un-j)
 =
 v_p((q+1)!).
\]
On the other hand,
\[
 v_p((d-1)!)
 =
 \sum_{\substack{1\le i<d\\p^a\nmid i}}v_p(i)
 +aq+v_p(q!).
\]
Subtracting $v_p(d!)=v_p((d-1)!)+v_p(d)$ therefore gives
\[
\begin{aligned}
 \min_n v_p\binom{mn}{d}
 &=
 a+v_p(q+1)-v_p(d).
\end{aligned}
\]
Hence
\[
 v_p(G_{m,r})
 =
 a_p+v_p(q_p+1)-v_p(d)
\]
for every $p\mid m$, while it is zero for $p\nmid m$. This proves
\eqref{eq:exact-gcd}.

Finally, write $g=\gcd(m,d)=\gcd(m,r-1)$. The exponent furnished by
\eqref{eq:exact-gcd} is at least $v_p(m)-v_p(g)$ for every $p\mid m$,
because $v_p(q_p+1)\ge0$. Therefore
\[
 \frac{m}{\gcd(m,r-1)}\mid G_{m,r}.
\]
\end{proof}

\begin{remark}[Sharpness criterion]
\label{rem:sharpness}
Theorem~\ref{thm:exact-gcd} gives a precise criterion for when the
universal divisor in Theorem~\ref{thm:div} is sharp. Namely,
\[
 G_{m,r}=\frac{m}{\gcd(m,r-1)}
\]
if and only if, for every prime $p\mid m$,
\[
 v_p\!\left(
 \left\lceil\frac{m+r-1}{p^{v_p(m)}}\right\rceil
 \right)
 =0.
\]
Thus the excess over the universal divisor is exactly the product of the
additional $p$-power factors contributed by these ceiling terms.
\end{remark}

\begin{lemma}[$2$-adic valuation of $\binom{2^s}{j}$]
\label{lem:2adic}
For every integer $s\ge1$ and every $0<j<2^s$,
\[
 v_2\left(\binom{2^s}{j}\right)=s-v_2(j).
\]
\end{lemma}

\begin{proof}
By Legendre's formula, $v_2(n!)=n-s_2(n)$, where $s_2(n)$ denotes the
number of $1$'s in the binary expansion of $n$. Thus
\[
 v_2\left(\binom{2^s}{j}\right)
 =-1+s_2(j)+s_2(2^s-j).
\]
If $t=v_2(j)$, then the binary expansions of $j$ and $2^s-j$ satisfy
\[
 s_2(j)+s_2(2^s-j)=s+1-t,
\]
which gives the stated formula.
\end{proof}

\begin{example}[Non-sharpness and the exact gcd]
\label{ex:nonsharp}
Take $m=2$ and $r=2^s-2$, with $s\ge2$. Then
\[
 d=m+r-1=2^s-1
\]
and Theorem~\ref{thm:decimation} gives
\[
 B_{2,r,k}=\binom{2^s}{2k-1}.
\]
Since every occurring exponent $2k-1$ is odd, Lemma~\ref{lem:2adic}
gives
\[
 v_2(B_{2,r,k})=s.
\]
Thus
\[
 \gcd_{1\le k\le K_{2,r}}B_{2,r,k}=2^s.
\]
Theorem~\ref{thm:div}, by contrast, gives only
\[
 \frac{2}{\gcd(2,2^s-3)}=2.
\]
This family therefore exhibits an unbounded gap, while
Theorem~\ref{thm:exact-gcd} recovers the exact value:
\[
 q_2=\left\lfloor\frac{2^s-2}{2}\right\rfloor=2^{s-1}-1,
\]
so
\[
 v_2(G_{2,r})
 =
 1+v_2(2^{s-1})-v_2(2^s-1)
 =s.
\]
\end{example}

\section{Row sums and Fourier multisection}
\label{sec:rows}

\begin{theorem}[Row-sum identity]
\label{thm:rowsum}
For every $m\ge2$ and $r\ge1$,
\[
 \sum_{k=1}^{K_{m,r}}B_{m,r,k}=m^{m+r-1}.
\]
\end{theorem}

\begin{proof}
Let
\[
 H(x)=(1+x+\cdots+x^{m-1})^{m+r}
\]
and let $\omega=e^{2\pi i/m}$. By Theorem~\ref{thm:decimation}, the desired
sum is the sum of coefficients of $H$ whose exponents are congruent to
$-1\pmod m$. The roots-of-unity filter gives
\[
 \sum_{k=1}^{K_{m,r}}B_{m,r,k}
 =
 \frac1m\sum_{j=0}^{m-1}\omega^jH(\omega^j).
\]
For $j=0$,
\[
 H(1)=m^{m+r}.
\]
For $1\le j\le m-1$,
\[
 1+\omega^j+\cdots+\omega^{j(m-1)}=0.
\]
Hence all nontrivial terms vanish and
\[
 \sum_{k=1}^{K_{m,r}}B_{m,r,k}
 =\frac1m m^{m+r}
 =m^{m+r-1}.
\]
\end{proof}

\begin{remark}
The row-sum formula says that the selected residue class carries exactly
one $m$-th of the total coefficient mass of
$(1+x+\cdots+x^{m-1})^{m+r}$. It is therefore a direct Fourier consequence
of the decimation theorem rather than an independent numerical
observation.
\end{remark}

\subsection{A Catalan specialization}

\begin{corollary}
\label{cor:catalan}
For $r=1$,
\[
 B_{m,1,1}
 =
 \binom{2m-1}{m}
 =
 \frac12\binom{2m}{m}
 =
 \frac{m+1}{2}C_m,
\]
where
\[
 C_m=\frac1{m+1}\binom{2m}{m}
\]
is the $m$-th Catalan number.
\end{corollary}

\begin{proof}
By Theorem~\ref{thm:decimation},
\[
 B_{m,1,1}
 =[x^{m-1}](1+x+\cdots+x^{m-1})^{m+1}.
\]
Since the target degree is $m-1$, the upper bounds on the individual
exponents do not bind, and the coefficient is
\[
 \binom{(m+1)+(m-1)-1}{m-1}
 =\binom{2m-1}{m-1}
 =\binom{2m-1}{m}.
\]
The remaining identities are standard.
\end{proof}

\section{Examples}

\begin{example}
For $(m,r)=(3,1)$, $d=3$ and $K_{3,1}=3$. We obtain
\[
 \binom{3n}{3}
 =
 10\binom{n}{3}
 +16\binom{n+1}{3}
 +\binom{n+2}{3}.
\]
The decimation representation is
\[
 (10,16,1)
 =
 \bigl([x^2],[x^5],[x^8]\bigr)(1+x+x^2)^4.
\]
The row sum is $27=3^3$, while the universal divisor is
$3/\gcd(3,0)=1$.
\end{example}

\begin{example}
For $(m,r)=(3,2)$, $d=4$ and $K_{3,2}=3$:
\[
 \binom{3n}{4}
 =
 15\binom{n}{4}
 +51\binom{n+1}{4}
 +15\binom{n+2}{4}.
\]
The vector $(15,51,15)$ is palindromic, as predicted by
$r\equiv2\pmod3$. Since
\[
 \frac{3}{\gcd(3,1)}=3,
\]
every coefficient is divisible by $3$, and the row sum is
$81=3^4$.
\end{example}

\begin{example}
For $(m,r)=(2,1)$,
\[
 \binom{2n}{2}
 =
 3\binom{n}{2}+\binom{n+1}{2}.
\]
The support has length $K_{2,1}=2$, the vector $(3,1)$ is not palindromic,
and the symmetry criterion correctly fails because $1\not\equiv0\pmod2$.
\end{example}

\section{Discussion of novelty and scope}

The coefficient identity \eqref{eq:main-decimation} should be distinguished
from the classical theory of multinomial coefficients. The coefficients of
$(1+x+\cdots+x^{m-1})^N$ are standard $m$-nomial coefficients. What is new
in the present formulation is their exact occurrence as the shifted-binomial
coordinates of $\binom{mn}{m+r-1}$, together with the consequences for the
coordinate array.

The universal divisibility theorem applies uniformly for all admissible
values of $m$ and $r$, rather than only in the cases where the coefficient
vector is palindromic.

The exact greatest common divisor
\[
 \gcd_{1\le k\le K_{m,r}}B_{m,r,k}
\]
is given by Theorem~\ref{thm:exact-gcd}. The formula shows explicitly why
the universal divisor in Theorem~\ref{thm:div} can fail to be sharp: the
additional factors are controlled by the $p$-adic valuations of the
quantities
\[
 \left\lceil\frac{m+r-1}{p^{v_p(m)}}\right\rceil,
 \qquad p\mid m.
\]
The family $m=2$, $r=2^s-2$ shows that these additional factors can be
arbitrarily large. Thus the exact row gcd is governed by local
$p$-adic information, but no longer represents an open problem in the
present setting. Natural further directions include inequalities such as
log-concavity for the decimated rows, and analogous coordinate identities
for other polynomial families admitting binomial-basis expansions.

\section{Conclusion}

We have studied the shifted-binomial coordinates of
\[
 \binom{mn}{m+r-1}.
\]
The central identity
\[
 B_{m,r,k}
 =
 [x^{mk-1}](1+x+\cdots+x^{m-1})^{m+r}
\]
identifies the coordinates with a residue-class decimation of a multinomial
coefficient sequence. From this representation we obtained the exact
support and positivity, and we proved the reflection criterion
$r\equiv2\pmod m$. We established the universal divisibility law
\[
 \frac{m}{\gcd(m,r-1)}\mid B_{m,r,k},
\]
and, more strongly, determined the exact row gcd by the explicit
$p$-adic formula in Theorem~\ref{thm:exact-gcd}. The formula also gives a
complete sharpness criterion for the universal divisor. Finally, a
roots-of-unity filter yields the exact row sum $m^{m+r-1}$.

The resulting framework separates the classical ingredients---binomial
bases, finite differences, multinomial coefficients, and roots-of-unity
filters---from the new structural identification that connects them.

\begin{thebibliography}{99}

\bibitem{CahenChabert1997}
P.-J. Cahen and J.-L. Chabert,
\emph{Integer-Valued Polynomials},
Mathematical Surveys and Monographs, Vol.~48,
American Mathematical Society, Providence, RI, 1997.

\bibitem{Comtet1974}
L. Comtet,
\emph{Advanced Combinatorics: The Art of Finite and Infinite Expansions},
D. Reidel Publishing Co., Dordrecht, 1974.

\bibitem{Polya1919}
G. P\'olya,
\"Uber ganzwertige Polynome in algebraischen Zahlk\"orpern,
\emph{J. Reine Angew. Math.} \textbf{149} (1919), 97--116.

\bibitem{Ostrowski1919}
A. Ostrowski,
\"Uber ganzwertige Polynome in algebraischen Zahlk\"orpern,
\emph{J. Reine Angew. Math.} \textbf{149} (1919), 117--124.
\bibitem{Riordan1968}
J. Riordan,
\emph{Combinatorial Identities},
John Wiley \& Sons, New York, 1968.

\bibitem{StanleyEC1}
R. P. Stanley,
\emph{Enumerative Combinatorics, Vol. 1},
2nd ed., Cambridge Studies in Advanced Mathematics, Vol.~49,
Cambridge University Press, Cambridge, 2011.

\end{thebibliography}

\end{document}
